Un circuito biohíbrido es un circuito formado por una neurona artificial (modelo) y una neurona viva (neurona real). En este, el modelo puede ser o bien un programa, ejecutándose en un ordenador, o un circuito eléctrico, que represente mediante sus componentes las partes de la célula. El circuito híbrido, puede ser unidireccional (con la interacción ocurriendo en una sola dirección) o bidireccional \cite{ReyesSnchez2023}. 
